\documentclass[11pt]{article}

\usepackage[T1]{fontenc}
\usepackage[margin=1in]{geometry}
\usepackage{amsmath,amssymb}
\usepackage[hidelinks]{hyperref}

\newcommand{\R}{\mathbb{R}}
\newcommand{\Z}{\mathbb{Z}}

\title{Literature-Implied Status of Shearlet Riesz-Basis Existence\\\large Source: arXiv:1404.1690}
\author{FA/Banach run \texttt{fa\_banach\_001}, lane 7}
\date{August 13, 2026}

\begin{document}
\maketitle

\section*{Status}

\texttt{literature\_implied\_answer}, scoped as: broad existence positive;
compact cone-adapted subclass open. This packet records a theorem-to-question
identification, not a new proof.

\section*{Source question}

Ma and Petersen study admissible compactly supported separable cone-adapted
shearlet systems in \cite{MaPetersen}. In Section 5, ``Beyond linear
independence,'' PDF page 18, immediately after their Lemma 5.2, they write:
\begin{quote}
The question about the existence of $\omega$-independent shearlet frames, or
equivalently shearlet Riesz bases is a delicate question for future work.
\end{quote}
Their Lemma 5.2 states, in particular, that a frame is a Riesz basis if and
only if it is $\omega$-independent.

\section*{Supporting theorem and identification}

F\"uhr and Maus construct what their abstract calls the first known shearlet
Riesz basis \cite{FuhrMaus}. Their Corollary 4.3, on PDF page 9, states that
for any rational $a>1$ and $b>0$ there exist $\psi\in L^2(\R^2)$ and
$\Lambda\subset\R^2$ such that
\[
 \left(D_{A_{a^k}}D_{S_{mb}}T_\lambda\psi\right)_
 {k,m\in\Z,\,\lambda\in\Lambda}
\]
is a Riesz basis of $L^2(\R^2)$. Hence it is also an
$\omega$-independent shearlet frame. This gives a positive answer to the
literal broad existence question in \cite{MaPetersen}.

The supporting paper does not cite Ma--Petersen or state that it is answering
their question. The relation is therefore an agent-identified implication,
not an explicit same-problem claim by the supporting authors.

\section*{Scope limitation}

The supporting construction is outside the exact class studied in the source
paper. It uses a wavelet-set generator, with $\widehat\psi$ proportional to
the characteristic function of a bounded frequency set, and the translation
set $\Lambda$ is a quasicrystal. It therefore does not construct a spatially
compactly supported, separable, cone-adapted shearlet Riesz basis with the
source paper's regular lattice sampling. Indeed, Führ--Maus Remark 4.5, also
on PDF page 9, only says that cone-adapted shearlet Riesz bases are expected.
The contextual compactly supported cone-adapted/lattice subproblem remains
open on the evidence reviewed here.

\section*{Search evidence and review recommendation}

The run indexes were searched using both arXiv ids and shorter keyword
combinations involving shearlets, Riesz bases, omega independence, cones, and
compact support. Local
primary-source comparison and bounded web/arXiv searches found the decisive
Führ--Maus construction but no later theorem closing the compactly supported
cone-adapted subclass. Review this only as a scoped literature-status packet;
do not count it as a new mathematical result or as a solution of that narrower
subproblem.

\raggedright
\begin{thebibliography}{9}

\bibitem{MaPetersen}
J. Ma and P. Petersen,
\emph{Linear independence of compactly supported separable shearlet systems},
arXiv:1404.1690.

\bibitem{FuhrMaus}
H. F\"uhr and Y. Maus,
\emph{Wavelet Riesz Bases Associated to Nonisotropic Dilations},
arXiv:1510.01832.

\end{thebibliography}

\end{document}
